\subsection*{A 26-modulus extremizer}
One further exact value is $m(26)=70$. Each of the 351 subsets of size 25 of the 27-element cap-60 pool already covers for some residues. The sole additional eligible modulus through cap 69 is 63, so every 26-set there contains a proper covering subset. This also appears in the companion calculation \cite{parallel}. At cap 70 add the one new modulus 70. An extremizer must therefore contain both 63 and 70. Of the 24-element parts from the cap-60 pool, 35 fail to cover. For each resulting 26-set a safe preliminary test discards it if any deletion already covers using at most seven 7-rows. On exactly seven selected rows the contraction above is valid, regardless of whether other rows were discarded. Omitting the preliminary test only enlarges the exhaustive search. It leaves 24 candidates. An unmatched singleton among 16 and 48 is also redundant.

Here is the unique passing set (the accompanying file gives its residues):
\[
 \{4,5,6,7,8,9,10,12,14,15,16,18,20,21,24,28,30,35,40,42,45,48,56,60,63,70\}.
\]
The following explains the final exact test, including why the two excess 7-rows cannot simply be discarded. After the binary contraction to 24, use a cycle of size 360 underneath the 7-digit. A hole of the ordinary rows can be moved to zero. Of the $7$ to $9$ top rows select seven distinct rows covering the seven lifts of that hole. Their projected residues must all be zero, and their upper digits can be named in their order. Enumerate all bucket and cofactor-residue choices for the remaining at most two rows. At each base point intersect the unions belonging to the seven buckets as before. If the ordinary rows leave no hole they can cover alone, which any such test permits. Keep the inclusion-maximal 360-bit masks up to translation. When all nine rows occur, twenty such masks suffice.

For each translate of each allowed mask, split its coordinates into the three highest ternary digits, writing $H_0,H_1,H_2$ on the base of 120. For each assignment of the ordinary 9-divisible rows (their projected moduli divide 120), let $V_j$ be the union covered by the rows with highest ternary digit $j$. The other rows have to cover the complement of the 120-bit mask
\[
                 (H_0\cup V_0)\cap(H_1\cup V_1)\cap(H_2\cup V_2).
\]
Try all allowed masks and fill each complement by direct residue enumeration. To reduce repetition one may translate any lower prime-power coordinate while fixing its highest digit, or relabel the highest ternary digit. Each such permutation maps every congruence class to a class of the same modulus. The full family of attainable 7-masks is invariant, and taking maximal masks up to translation preserves existential coverability. Thus these normalizations can be applied to the ordinary rows before choosing the 7-mask. Run the decision again after every single deletion. The exhaustive list contains only the displayed 26-set. A separate decision code constructs 7-masks from unlabeled partitions with a singleton bucket, projects the highest binary digit onto 180, and verifies that the displayed set covers and every deletion fails. The explicit positive residues were also checked directly over period 5040. Lean 4.33.1 checks that the 26 moduli are distinct, their maximum is 70, and the residue assignment covers every integer. The irreducibility and optimality arguments use the finite algorithms.
